\documentclass[10pt,a4paper]{amsart}
\usepackage[margin=19mm]{geometry}
\usepackage{amsmath,amssymb,mathtools}
\usepackage[scr=rsfs,cal=euler]{mathalfa}
\usepackage{xfrac}
\usepackage[unicode,hidelinks,bookmarks=false]{hyperref}
\newcommand{\ie}{i.e.\ }
\newcommand{\cf}{cf.\ }
\newcommand{\resp}{respectively\ }
\newcommand{\Subsection}{\subsection}
\providecommand{\nobreakdash}{\nobreak\hbox{-}}
\setlength{\emergencystretch}{2em}
\multlinegap0pt
\usepackage{bm}
\usepackage{bbm}
\usepackage{dsfont}
\usepackage{xspace}
\usepackage{cancel}
\usepackage{mathtools}
\usepackage{cleveref}
\usepackage[shortlabels]{enumitem}
\usepackage{amsthm}
\makeatletter
\@ifundefined{THM}{\newtheorem{THM}{Thm}}{}
\@ifundefined{LEM}{\newtheorem{LEM}[THM]{Lemma}}{}
\@ifundefined{PROP}{\newtheorem{PROP}[THM]{Proposition}}{}
\makeatother
\makeatletter
\def\G{{\Gamma}}
\DeclareMathOperator{\PMap}{PMap}
\DeclareMathOperator{\PMapc}{PMap_{\it c}}
\newcommand{\cE}{\mathcal{E}}
\newcommand{\cP}{\mathcal{P}}
\newcommand{\cV}{\mathcal{V}}
\expandafter\ifx\csname customcor\endcsname\relax
\newenvironment{customcor}[1]
  {\renewcommand\theinnercustomcor{#1}\innercustomcor}
  {\endinnercustomcor}
\fi
\expandafter\ifx\csname customthm\endcsname\relax
\newenvironment{customthm}[1]
  {\renewcommand\theinnercustomthm{#1}\innercustomthm}
  {\endinnercustomthm}
\fi
\DeclareMathOperator{\id}{id}
\makeatother
\title[Generating Sets and Algebraic Properties of Pure Mapping Cla]{Generating Sets and Algebraic Properties of Pure Mapping Class Groups of Infinite Graphs}
\author{George Domat, Hannah Hoganson, Sanghoon Kwak}
\date{Source: arXiv:2309.07885v1 (CC BY 4.0)}
\begin{document}
\maketitle

\noindent\emph{Source and licence.} Generating Sets and Algebraic Properties of Pure Mapping Class Groups of Infinite Graphs. Version arXiv:2309.07885v1 (CC BY 4.0), distributed under the Creative Commons Attribution 4.0 International licence, as recorded in the arXiv licence metadata for this exact version and shown on the abstract page https://arxiv.org/abs/2309.07885v1. Full rights notice: licenses/arxiv-2309.07885-CC-BY-4.0.txt.\par\smallskip

\noindent\textit{Editorial scope.} Counted unit: the Prop whose source label is \texttt{prop:coregraphcase} in the article (source0.tex lines 814--816, source label \texttt{prop:coregraphcase}), delivered together with the article's own complete original proof of it (source0.tex lines 818--842, one \texttt{proof} environment of 25 lines). No mathematics is written, repaired or re-derived: statement and proof are the article's own text line for line. Required context retained uncounted: lem:zeroflux (lines 753--763). Every definition, display and label of those lines is retained unchanged. Omitted from this package and not redistributed: 1--752, 764--813, 817--817, 843--1550. Nothing inside a retained excerpt is omitted or altered: every retained body is delivered exactly as the article's lines are written, so no omission receipt of any kind accompanies this package. Citations: the retained excerpts carry no citation command at all, so no bibliography entry of the article is transcribed and no reference is redistributed. Reference adaptation: none was needed and none was made; every reference inside the delivered unit resolves inside it, so no unresolved reference or citation is produced. Printed slips kept literal: no printed word, symbol, hypothesis or number of the counted statement or of its proof was altered. Layout and class: the article's own class is \texttt{amsart} with packages including inputenc, amsmath, amssymb, amsthm, enumitem, xcolor, wasysym, booktabs, makecell, tabularx, colortbl, multirow, adjustbox, caption; the wrapper is the lane's pinned \texttt{amsart} editorial wrapper at 10pt on A4, which restates the theorem-like environments the retained text needs and adds the article's own macro definitions that the retained text needs (\texttt{\textbackslash G}, \texttt{\textbackslash PMap}, \texttt{\textbackslash PMapc}, \texttt{\textbackslash cE}, \texttt{\textbackslash cP}, \texttt{\textbackslash cV}, \texttt{\textbackslash id}), with extra wrapper packages (bm, bbm, dsfont, xspace, cancel, mathtools, cleveref, enumitem). The delivered numbering is the unit's own. Assets: no retained span contains a figure environment, a graphics-inclusion command, a PDF-page inclusion, a table or a TikZ picture; no image file is copied into this package, the package declares no asset dependency and no image file is redistributed with it. Licence: version arXiv:2309.07885v1 of this article is distributed under the Creative Commons Attribution 4.0 International licence, as recorded in the arXiv licence metadata for this exact version and shown on the abstract page https://arxiv.org/abs/2309.07885v1. A full rights notice is delivered at licenses/arxiv-2309.07885-CC-BY-4.0.txt. Characters: every retained excerpt is pure ASCII, so the delivered engine, XeLaTeX with its default Latin Modern text font, reports no missing character.
\begin{LEM} \label{lem:zeroflux}
    Let $\cE \subset E(\G)$ be a nonempty proper clopen subset and $f \in \PMap(\G)$. If $\Phi_{\cE}(f) = 0$, then given any compact $K \subset \G$, there exists $\psi \in \PMap(\G)$ such that
    \begin{enumerate}[label=(\roman*)]
        \item \textrm{\bf (Compact approximation)} 
            $\psi f^{-1} \in \cV_{K}$,
        \item \textrm{\bf (Truncation)} there exist disjoint subgraphs $\G_{\cE}$, and $\G_{\cE^{c}}$ of $\G$ with end spaces $\cE$ and $\cE^{c}$ respectively, such that 
            $\psi\vert_{\G_{\cE}} = \id$ and $\psi\vert_{\G_{\cE^{c}}} = f\vert_{\G_{\cE^{c}}}$, and
        \item \textrm{\bf (Same flux)}
            $\Phi_{\eta}(\psi) = \Phi_{\eta}(f)$ for every clopen subset $\eta \subset E(\G) \setminus \cE$.
    \end{enumerate}
\end{LEM}

\begin{PROP}\label{prop:coregraphcase}
    Let $\G$ be a locally finite, infinite graph with $E(\G) = E_{\ell}(\G)$, $|E(\G)| \geq 2$, and $f \in \PMap(\G)$. If $\Phi_{{\cE}}(f) = 0$ for every clopen subset $\cE$ of $E(\G)$, then $f \in \overline{\PMap_c(\G)}$.
\end{PROP}

\begin{proof}
    Assume $f \in \PMap(\G)$ has $\Phi_{\cE}(f) = 0$ for every nonempty proper clopen subset $\cE$ of the end space $E(\G)$. Given any compact $K \subset \G$ we will find $\psi \in \PMapc(\G)$ such that $\psi f^{-1} \in \cV_{K}$. 

    Without loss of generality we may enlarge $K$ so that it is connected, has at least two complementary components, and every complementary component of $K$ is infinite. Then the complement of $K$ induces a partition of the ends. Write 
    \[
    \cP_K = \cE_1 \sqcup \ldots \sqcup \cE_n
    \]
    for this partition.

    Apply \Cref{lem:zeroflux} to $f$ using $\cE_{1}$ to obtain $\psi_{1}$. Note that by (iii) we still have $\Phi_{\cE_{2}}(\psi_{1}) = \Phi_{\cE_2}(f) = 0$. Thus we can apply the lemma again to $\psi_{1}$ using $\cE_{2}$ to obtain a $\psi_{2}$. Continue this process recursively to obtain $\psi_{n}$. 

    Now, by (i) of \Cref{lem:zeroflux}, there exist $v_1,\ldots,v_n \in \cV_K$ such that
    \begin{align*}
        \psi_{i} &=
        \begin{cases}
            v_i\psi_{i-1} & \text{for $1 < i \le n$,}\\
            v_1f & \text{for $i=1$}.
        \end{cases}
    \end{align*}
    Putting these together gives $\psi_{n}f^{-1} = v_{n}v_{n-1}\cdots v_{1} \in \cV_{K}$ as $\cV_{K}$ is a subgroup.

    It remains to check that $\psi_{n} \in \PMapc(\G)$. However, by (ii), we have that $\psi_{n}$ is equal to the identity on $\bigcup_{i=1}^{n}\G_{\cE_{i}}$. This exactly covers all of the ends of $\G$ as $\cP_{K}$ was a partition of the ends. Therefore we see that $\psi_{n}$ is supported on $\overline{\bigcap_{i=1}^{n} \G \setminus \G_{\cE_{i}}}$, a compact set. Taking $\psi = \psi_{n}$ gives the desired compact approximation of $f$. 

    Finally, since the $K$ above was taken to be arbitrary, starting with a compact exhaustion of $\G$ we can apply the above to obtain a sequence of compactly supported maps that converge to $f$. 
\end{proof}

\end{document}
